\begin{pastexam}[label={px:ust-115-ee-ds-29}]{臺聯大電機類 115 學年度 第 29 題}
（單選；2 分）

Answer \textbf{Questions 28-30} according to the following C code.
\begin{minted}[linenos, numbersep=1.5mm, xleftmargin=2em, fontsize=\small]{c}
#include <stdio.h>
#include <stdlib.h>
int counter=0;
void max_heapify(int A[],int n,int i) {
    int largest=i;
    int left=2*i+1;
    int right=2*i+2;
    counter++;
    if (left<n && A[left]>A[i])
        largest=left;
    if (right<n && A[right]>A[largest])
        largest=right;
    if (largest!=i) {
        int temp=A[i];
        A[i]=A[largest];
        A[largest]=temp;
        max_heapify(A,n,largest);
    }
}
void build_heap(int A[],int n) {
    for (int i = (n/2)-1;i>=0;i--)
        max_heapify(A,n,i);
}
int main(){
    int A[]={4,1,3,2,16,9,10,14,8,7};
    int n = sizeof(A)/sizeof(A[0]);
    build_heap(A,n);
    for (int i=0;i<n;++i)
        printf("%d,",A[i]);
    printf("\n");
    printf("counter:%d\n",counter);
    return 0;
}
\end{minted}

What does the code at \textbf{Line 31} print? Counter：\underline{\hspace{2em}}

(A)~6 \qquad (B)~10 \qquad (C)~12 \qquad (D)~14 \qquad (E)~16
\end{pastexam}

\begin{pastsoln}
答案是 (C)。

\texttt{counter} 在 \texttt{max\_heapify} 的開頭加一，所以它等於 \texttt{max\_heapify} 被呼叫的
總次數，包括遞迴呼叫。 \texttt{build\_heap} 對 $i = 4, 3, 2, 1, 0$ 各呼叫一次，每次往下沉幾層，
就多幾次遞迴呼叫（最後一次是發現不必交換、直接結束的那一次）。 沿用上一題（歷屆試題~\ref{px:ust-115-ee-ds-28}）
的追蹤：
\begin{center}
\begin{tabular}{lll}
\toprule
$i$ & 呼叫的足標 & 次數 \\
\midrule
4 & 4 & 1 \\
3 & 3, 7 & 2 \\
2 & 2, 6 & 2 \\
1 & 1, 4, 9 & 3 \\
0 & 0, 1, 3, 8 & 4 \\
\midrule
合計 & & 12 \\
\bottomrule
\end{tabular}
\end{center}
所以印出 \texttt{counter:12}。

其他選項：(A) $6$ 只比外層的 $5$ 次多一次，漏算了大部分遞迴呼叫；(B) $10$ 是 $n$，不是呼叫次數；
(D)(E) 都比實際多。 另一種數法：交換共 $7$ 次，外層呼叫 $5$ 次，每次交換恰好多一次遞迴呼叫，$5 + 7 = 12$，和表格一致。
\end{pastsoln}
